\rok{Јануар 2024.}{D}

Други практични тест из Алгоритама и структура података

Други практични тест одржан 09.01.2024.

\zadatak{1}{Конверзија матрице у листу суседства}
Написати алгоритам који ће вршити директну конверзију матрице суседства у листу суседства. Алгоритам као улазни параметар треба да прихвати матрицу, а треба да резултује креираним динамичким низом чворова.

\textbf{Додатне напомене за решавање задатка:}
\begin{itemize}
	\item У template-у је закоментарисана метода \texttt{ConvertAdjacencyMatrixToAdjacencyList} која треба да садржи имплементацију конверзије матрице у листу суседства.
	\item У Main методи обезбедити начин за тестирање и проверу нове функционалности.
\end{itemize}



\zadatak{2}{Пребројавање понављајућих вредности у низу}
Написати алгоритам временске комплексности $O(n)$ који ће претражити низ целих бројева и пребројати све вредности које се понављају.

\textbf{Додатне напомене за решавање задатка:}
\begin{itemize}
	\item Обавезно је користити Hash табелу која је дата у шаблону.
	\item Имплементацију писати у оквиру закоментарисане методе:
	\begin{Verbatim}[fontsize=\footnotesize, numbers=left, xleftmargin=10mm]
public static int CountRepeatedDistinctValues(int[] arr)
	\end{Verbatim}
\end{itemize}

\textbf{Примери:}
\begin{itemize}
	\item \textbf{Улаз 1:} \texttt{[5, 6, 3, 4, 3, 6, 4]}
	\item \textbf{Излаз 1:} 3
	\item \textbf{Улаз 2:} \texttt{[1, 2, 3, 4, 5, 6]}
	\item \textbf{Излаз 2:} 0
	\item \textbf{Улаз 3:} \texttt{[1, 1, 1, 1, 1]}
	\item \textbf{Излаз 3:} 1
\end{itemize}



\zadatak{3}{Исписивање чворова на позицији листова}
У оквиру стандардне структуре класе \texttt{AVL} написати имплементацију методе која ће омогућити исписивање свих чворова који се налазе на позицији листова.

\textbf{Додатне напомене за решавање задатка:}
\begin{itemize}
	\item У оквиру класе \texttt{AVL} креирати јавну методу са потписом:
	\begin{Verbatim}[fontsize=\footnotesize, numbers=left, xleftmargin=10mm]
public void printLeaves()
	\end{Verbatim}
	\item Из претходно креиране јавне методе треба да се позива приватна метода са потписом:
	\begin{Verbatim}[fontsize=\footnotesize, numbers=left, xleftmargin=10mm]
private void printLeavesAVL(Node curr)
	\end{Verbatim}
	\item Приватна метода треба да садржи имплементацију алгоритма којим ће се омогућити пролазак кроз стабло од корена ка листовима и њихово пребројавање.
	\item Листом се класификује сваки чвор који нема ниједно дете.
	\item Алгоритам \textbf{ОБАВЕЗНО} писати у оквиру класе \texttt{AVL}.
\end{itemize}

\textbf{Примери:}

\begin{center}
\begin{minipage}{\textwidth}
\centering
\begin{forest}
for tree={
	circle,
	draw,
	thick,
	fill=gray!20,
	minimum size=8mm,
	inner sep=1pt,
	s sep=8mm,
	l sep=12mm,
	edge={-, thick},
	font=\small
}
[8
	[4
		[2
			[, phantom]
			[3]
		]
		[5]
	]
	[9
		[, phantom]
		[11]
	]
]
\end{forest}

\vspace{0.3cm}
\textbf{Слика 1.} Пример BST-а.
\end{minipage}
\end{center}

\begin{itemize}
	\item \textbf{Улаз 1:} BST са слике 1
	\item \textbf{Излаз 1:} 3, 5, 11
\end{itemize}

\begin{center}
\begin{minipage}{\textwidth}
\centering
\begin{forest}
for tree={
	circle,
	draw,
	thick,
	fill=gray!20,
	minimum size=8mm,
	inner sep=1pt,
	s sep=8mm,
	l sep=12mm,
	edge={-, thick},
	font=\small
}
[10
	[4
		[2
			[3]
			[, phantom]
		]
		[6
			[5]
			[9]
		]
	]
	[20
		[12
			[, phantom]
			[14]
		]
		[29
			[27]
			[, phantom]
		]
	]
]
\end{forest}

\vspace{0.3cm}
\textbf{Слика 2.} Пример BST-а.
\end{minipage}
\end{center}

\begin{itemize}
	\item \textbf{Улаз 2:} BST са слике 2
	\item \textbf{Излаз 2:} 3, 5, 9, 14, 27
\end{itemize}



\sablon{Шаблон другог теста јануар 2024. група D}{2023/Januar/GrupaD/AISP2023_Test2_GrupaD.linq}
